\documentclass[12pt,a4paper]{report}

% -------------------------------------------------
% Packages
% -------------------------------------------------
\usepackage[a4paper,margin=1in]{geometry}
\usepackage{amsmath,amssymb,amsthm}
\usepackage{physics}
\usepackage{graphicx}
\usepackage{float}
\usepackage{caption}
\usepackage{hyperref}
\usepackage{fancyhdr}
\usepackage{setspace}
\usepackage{titlesec}

% -------------------------------------------------
% Custom Commands
% -------------------------------------------------
\newcommand{\vecspace}[1]{\mathbb{#1}}

% -------------------------------------------------
% Header & Footer
% -------------------------------------------------
\pagestyle{fancy}
\fancyhf{}
\fancyhead[L]{Quantum Mechanics Report}
\fancyhead[R]{\leftmark}
\fancyfoot[C]{\thepage}

\hypersetup{pdfborder=0 0 0}

\setlength{\headheight}{15pt}

% -------------------------------------------------
% Spacing
% -------------------------------------------------
\onehalfspacing

% -------------------------------------------------
% Title Information
% -------------------------------------------------
\title{
    \vspace{3cm}
    \Huge\textbf{QUANTUM MECHANICS}\\[1cm]
    \Large Mid-Term Report
}

\author{
    \textbf{Daksh Maahor} \\[0.3cm]
    25B0974 \\[0.3cm]
    Computer Science and Engineering \\[0.3cm]
    Mentor: Sourav Sreenath \\[0.3cm]
    Indian Institute of Technology Bombay
}

\date{\today}

% -------------------------------------------------
% Document
% -------------------------------------------------
\begin{document}

% -------------------------------------------------
% Cover Page
% -------------------------------------------------
\begin{titlepage}
\centering

\vspace*{2cm}

{\Large SUMMER OF SCIENCE 2026\\[0.2cm]
MID-TERM REPORT}

\vspace{1.5cm}
\rule{0.8\textwidth}{0.5pt}

\vspace{1cm}

{\Huge\bfseries QUANTUM MECHANICS}

\vspace{1cm}

\rule{0.8\textwidth}{0.5pt}

\vfill

{\Large
\textbf{Daksh Maahor}\\
25B0974\\
Computer Science and Engineering\\
1st Year (2026)\\[0.5cm]

Mentor: Sourav Sreenath\\
Indian Institute of Technology Bombay
}

\vfill

\end{titlepage}

% -------------------------------------------------
% Table of Contents
% -------------------------------------------------
\tableofcontents
\clearpage

% -------------------------------------------------
% Chapter 1
% -------------------------------------------------
\chapter{Mathematical Foundation}

\section{Introduction to Linear Vector Spaces}

A Linear Vector Space $\vecspace{V}$ is a mathematical structure that allows for the addition of vectors and multiplication by scalars. It contains a collection of objects called vectors, which have:

\begin{itemize}
    \item A definite rule for forming a vector sum $\ket{V} + \ket{W}$.
    \item A definite rule for forming a scalar product $c\ket{V}$, where $c$ is a complex number.
\end{itemize}
and the following features:
\begin{itemize}
    \item Closure under addition: If $\ket{V}$ and $\ket{W}$ are in $\vecspace{V}$, then $\ket{V} + \ket{W}$ is also in $\vecspace{V}$.
    \item Distributivity of scalar multiplication: $c(\ket{V} + \ket{W}) = c\ket{V} + c\ket{W}$.
    \item Distributivity of scalar multiplication with respect to field addition: $(c + d)\ket{V} = c\ket{V} + d\ket{V}$, where $c$ and $d$ are scalars.
    \item Commutativity of vector addition: $\ket{V} + \ket{W} = \ket{W} + \ket{V}$.
    \item Associativity of vector addition: $(\ket{V} + \ket{W}) + \ket{Z} = \ket{V} + (\ket{W} + \ket{Z})$.
    \item Existence of an additive identity: There exists a vector $\ket{0}$ such that $\ket{V} + \ket{0} = \ket{V}$.
    \item Existence of additive inverse: For every vector $\ket{V}$, there exists a vector $-\ket{V}$ such that $\ket{V} + \ket{-V} = \ket{0}$.
\end{itemize}

\subsection{Inner Product Spaces}

The inner product is a mathematical operation that takes two vectors and produces a scalar. It is denoted as $\bra{V}\ket{W}$ and had the following properties:
\begin{itemize}
    \item Conjugate symmetry: $\bra{V}\ket{W} = \bra{W}\ket{V}^*$.
    \item Linearity in the second argument: $\bra{V}\ket{aW + bX} = a\bra{V}\ket{W} + b\bra{V}\ket{X}$.
    \item Conjugate Linearity in the first argument: $\bra{cW + dZ}\ket{V} = c^*\bra{W}\ket{V} + d^*\bra{Z}\ket{V}$.
    \item Positive definiteness: $\bra{V}\ket{V} \geq 0$ and $\bra{V}\ket{V} = 0$ if and only if $\ket{V} = \ket{0}$.
\end{itemize}

Two vectors $\ket{V}$ and $\ket{W}$ are said to be orthogonal if their inner product is zero, i.e., $\bra{V}\ket{W} = 0$.

Also we define the norm of a vector $\ket{V}$ as $\|V\| = \sqrt{\bra{V}\ket{V}}$. A vector with unit norm is called a unit vector.
Two important inequalities related to inner product spaces are:

\begin{itemize}
    \item Cauchy-Schwarz Inequality: $|\bra{V}\ket{W}| \leq \|V\| \|W\|$.
    \item Triangle Inequality: $\|V + W\| \leq \|V\| + \|W\|$.
\end{itemize}

\subsection{Basis Vectors}

A set of vectors $\{\ket{e_1}, \ket{e_2}, \ldots, \ket{e_n}\}$ is called a basis for a vector space $\vecspace{V}$ if every vector in $\vecspace{V}$ can be expressed as a linear combination of the basis vectors. For example, if $\vecspace{V}^3$ is a 3-dimensional vector space, then any vector $\ket{V} \in \vecspace{V}^3$ can be written as:
\begin{equation*}
    \ket{V} = v_1\ket{e_1} + v_2\ket{e_2} + v_3\ket{e_3}
\end{equation*}
where $v_1$, $v_2$, and $v_3$ are complex numbers called the components of the vector $\ket{V}$ in the basis $\{\ket{e_1}, \ket{e_2}, \ket{e_3}\}$.

A basis is called orthonormal if the basis vectors are orthogonal to each other and each basis vector has a unit norm. In an orthonormal basis, the inner product of two basis vectors is given by:
\begin{equation*}
    \bra{i}\ket{j} = \delta_{ij}
\end{equation*}
where $\delta_{ij}$ is the Kronecker delta symbol, defined as:
\begin{equation*}
    \delta_{ij} = \begin{cases}
        1 & \text{if } i = j \\
        0 & \text{if } i \neq j
    \end{cases}
\end{equation*}

The component of a vector $\ket{V}$ in the direction of a basis vector $\ket{j}$ can be calculated using the inner product:
\begin{equation*}
    v_j = \bra{j}\ket{V}
\end{equation*}

Thus, in any basis, we can express a vector $\ket{V}$ as:

\begin{equation*}
    \ket{V} = \sum_{j} \ket{j}\bra{j}\ket{V}
\end{equation*}

Consider two vectors $\ket{V}$ and $\ket{W}$ expanded in a basis as:

\begin{equation*}
    \ket{V} = \sum_{i} v_i \ket{i}
\end{equation*}

\begin{equation*}
    \ket{W} = \sum_{j} w_j \ket{j}
\end{equation*}

The inner product can be expanded as:

\begin{equation*}
    \bra{V}\ket{W} = \left( \sum_{i} v_i^* \bra{i} \right) \left( \sum_{j} w_j \ket{j} \right) = \sum_{i,j} v_i^* w_j \bra{i}\ket{j}
\end{equation*}

Now if the basis was orthonormal, then $\bra{i}\ket{j} = \delta_{ij}$, and the inner product simplifies to:

\begin{equation*}
    \bra{V}\ket{W} = \sum_{i} v_i^* w_i
\end{equation*}

From now on, we will assume that all bases are orthonormal unless otherwise specified.

\subsection{Matrix Representation of Vectors}

The basis ket $\ket{i}$ is represented as a column vector with a 1 in the $i$-th position and 0s elsewhere.

In a finite-dimensional vector space, we can represent vectors as column matrices. For example, if $\vecspace{V}^3$ is a 3-dimensional vector space, a vector $\ket{V} \in \vecspace{V}^3$ can be represented as:
\begin{equation*}
    \ket{V} = v_1 \ket{1} + v_2 \ket{2} + v_3 \ket{3} = v_1 \begin{bmatrix}
        1 \\
        0 \\
        0
    \end{bmatrix} + v_2 \begin{bmatrix}
        0 \\
        1 \\
        0
    \end{bmatrix} + v_3 \begin{bmatrix}
        0 \\
        0 \\
        1
    \end{bmatrix} = \begin{bmatrix}
        v_1 \\
        v_2 \\
        v_3
    \end{bmatrix}
\end{equation*}

\subsection{Adjoint of a Vector}

The adjoint of a vector $\ket{V}$, denoted as $\bra{V}$, is the conjugate transpose of the vector. If $\ket{V}$ is represented as a column matrix, then $\bra{V}$ is represented as a row matrix:

\begin{equation*}
    \bra{V} = \begin{bmatrix}
        v_1^* & v_2^* & v_3^*
    \end{bmatrix}
\end{equation*}

With that representation, the inner product of two n-dimensional vectors $\ket{V}$ and $\ket{W}$ can be calculated as:
\begin{equation*}
    \bra{V}\ket{W} = \begin{bmatrix}
        v_1^* & v_2^* & v_3^* & \ldots & v_n^*
    \end{bmatrix}
    \begin{bmatrix}
        w_1 \\
        w_2 \\
        \vdots \\
        w_n
    \end{bmatrix}
    = v_1^*w_1 + v_2^*w_2 + \cdots + v_n^*w_n
\end{equation*}

Some properties of the adjoint operation are:
\begin{itemize}
    \item $(\ket{V} + \ket{W})^\dagger = \bra{V} + \bra{W}$.
    \item $(c\ket{V})^\dagger = \bra{V} c^*$, where $c$ is a complex number.
\end{itemize}

\subsection{Vector Subspaces}

A subspace $\vecspace{W}$ of a vector space $\vecspace{V}$ is a subset of $\vecspace{V}$ that is itself a vector space under the same operations of addition and scalar multiplication. For $\vecspace{W}$ to be a subspace of $\vecspace{V}$, it must satisfy the following conditions:
\begin{itemize}
    \item The zero vector $\ket{0}$ of $\vecspace{V}$ is in $\vecspace{W}$.
    \item If $\ket{V}$ and $\ket{W}$ are in $\vecspace{W}$, then $\ket{V} + \ket{W}$ is also in $\vecspace{W}$ (closure under addition).
    \item If $\ket{V}$ is in $\vecspace{W}$ and $c$ is a scalar, then $c\ket{V}$ is also in $\vecspace{W}$ (closure under scalar multiplication).
\end{itemize}
A vector subspace of dimension $n$ is denoted as $\vecspace{V}^n$. Given two vector subspaces $\vecspace{V}^n$ and $\vecspace{W}^m$, we can form their direct sum $\vecspace{V}^n \oplus \vecspace{W}^m$, which is a vector space of dimension $n + m$. This sum is defined as the set of all vectors that can be written as $\ket{V} + \ket{W}$, where $\ket{V} \in \vecspace{V}^n$ and $\ket{W} \in \vecspace{W}^m$.
\section{Linear Operators}
An operator $\hat{\Omega}$ is an instruction to transform any vector $\ket{V}$ to another vector $\ket{V'}$.
\begin{equation*}
    \hat{\Omega} \ket{V} = \ket{V'}
\end{equation*}
Operators can also act on bras in similar way:
\begin{equation*}
    \bra{V} \hat{\Omega} = \bra{V'}
\end{equation*}
An operator is called a linear operator if it satisfies the following properties:
\begin{itemize}
    \item Linearity in kets: $\hat{\Omega}(a\ket{V} + b\ket{W}) = a\hat{\Omega}\ket{V} + b\hat{\Omega}\ket{W}$, where $a$ and $b$ are scalars.
    \item Linearity in bras: $(\bra{V}a + \bra{W}b) \hat{\Omega} = \bra{V}\hat{\Omega} a + \bra{W} \hat{\Omega} b$, where $a$ and $b$ are scalars.
\end{itemize}

\subsection{Some useful Linear Operators}

\begin{itemize}

\item Identity Operator

The Identity operator $\hat{I}$ is the operator which maps every vector to itself:
\begin{align*}
    \hat{I} \ket{V} = \ket{V} \\
    \bra{V} \hat{I} = \bra{V}
\end{align*}

\item Rotation Operator

The operator $\hat{R}(\frac{\pi}{2}i)$ is the operator which rotates every vector by $\frac{\pi}{2}$ radians about the $i$-th axis. For example, in a 3-dimensional vector space, the operator $\hat{R}(\frac{\pi}{2}z)$ rotates every vector by $\frac{\pi}{2}$ radians about the $z$-axis.
Its effect on the basis vectors is given by:
\begin{align*}
    \hat{R}(\frac{\pi}{2}z) \ket{1} = \ket{2} \\
    \hat{R}(\frac{\pi}{2}z) \ket{2} = -\ket{1} \\
    \hat{R}(\frac{\pi}{2}z) \ket{3} = \ket{3}
\end{align*}

More generally, the operator $\hat{R}(\vec{\theta})$ represents a rotation by an angle $|\vec{\theta}|$ about the axis defined by the vector $\hat{\theta} = \frac{\vec{\theta}}{|\vec{\theta}|}$.

\item Projection Operator

Consider a vector $\ket{V}$. In a given basis, we can express $\ket{V}$ as:
\begin{equation*}
    \ket{V} = \sum_{i} \ket{i}v_i
\end{equation*}

Here, $v_i$ is the component of $\ket{V}$ in the direction of the basis vector $\ket{i}$, which can be written as $v_i = \bra{i}\ket{V}$. Thus,
\begin{equation*}
    \ket{V} = \sum_{i} \ket{i}\bra{i}\ket{V} = \left(\sum_{i} \ket{i}\bra{i} \right)\ket{V} = \hat{I} \ket{V}
\end{equation*}

Now, since this is true for any vector $\ket{V}$, we can write:
\begin{equation*}
    \hat{I} = \sum_{i} \ket{i}\bra{i} = \sum_{i} \hat{P}_i
\end{equation*}

Here, the operator $\hat{P}_i = \ket{i}\bra{i}$ is called the projection operator onto the basis vector $\ket{i}$. The projection operator $\hat{P}_i$ has the following properties:
\begin{itemize}
    \item $\hat{P}_i^2 = \hat{P}_i$
    \item $\hat{P}_i \hat{P}_j = \delta_{ij}$
    \item $\hat{P}_i \ket{V} = v_i \ket{i}$, which is the projection of $\ket{V}$ onto the basis vector $\ket{i}$.
\end{itemize}

\end{itemize}

\subsection{Matrix Representation of Operators}

For a linear operator $\hat{\Omega}$, given a basis $\ket{1}, \ket{2}, \ldots, \ket{n}$ for a vector space $\vecspace{V}^n$, and for any vector $\ket{V} \in \vecspace{V}^n$, we can express $\ket{V}$ as:
\begin{equation*}
    \ket{V} = \sum_{i} v_i \ket{i}
\end{equation*}
Then the action of the operator $\hat{\Omega}$ on $\ket{V}$ can be expressed as:
\begin{equation*}
    \hat{\Omega} \ket{V} = \hat{\Omega} \left( \sum_{i} v_i \ket{i} \right) = \sum_{i} v_i \hat{\Omega} \ket{i} = \sum_{i} v_i \ket{i'}
\end{equation*}
where $\ket{i'} = \hat{\Omega} \ket{i}$.

The operator $\hat{\Omega}$ can be represented as a matrix $\Omega$ in the basis $\ket{1}, \ket{2}, \ldots, \ket{n}$, where the $j$-th column of $\Omega$ is given by the vector $\ket{j'} = \hat{\Omega} \ket{j}$ expressed in the basis. Thus, we can write:
\begin{equation*}
    \Omega_{ji} = \bra{j}\hat{\Omega}\ket{i}
\end{equation*}

Using the fact that $\ket{i'} = \hat{\Omega} \ket{i}$, we can also write the matrix elements as:
\begin{equation*}
    \Omega = \begin{bmatrix}
        \bra{1}\hat{\Omega}\ket{1} & \bra{1}\hat{\Omega}\ket{2} & \cdots & \bra{1}\hat{\Omega}\ket{n} \\
        \bra{2}\hat{\Omega}\ket{1} & \bra{2}\hat{\Omega}\ket{2} & \cdots & \bra{2}\hat{\Omega}\ket{n} \\
        \vdots & \vdots & \ddots & \vdots \\
        \bra{n}\hat{\Omega}\ket{1} & \bra{n}\hat{\Omega}\ket{2} & \cdots & \bra{n}\hat{\Omega}\ket{n}
    \end{bmatrix}
\end{equation*}

This is the matrix representation of the operator $\hat{\Omega}$ in the basis $\ket{1}, \ket{2}, \ldots, \ket{n}$. In this representation, the action of the operator $\hat{\Omega}$ on a vector $\ket{V}$ can be calculated using matrix multiplication:
\begin{equation*}
    \ket{V'} = \Omega \ket{V}
\end{equation*}
\begin{equation*}
    \begin{bmatrix}
        v'_1 \\
        v'_2 \\
        \vdots \\
        v'_n
    \end{bmatrix}
    =
    \begin{bmatrix}
        \bra{1}\hat{\Omega}\ket{1} & \bra{1}\hat{\Omega}\ket{2} & \cdots & \bra{1}\hat{\Omega}\ket{n} \\
        \bra{2}\hat{\Omega}\ket{1} & \bra{2}\hat{\Omega}\ket{2} & \cdots & \bra{2}\hat{\Omega}\ket{n} \\
        \vdots & \vdots & \ddots & \vdots \\
        \bra{n}\hat{\Omega}\ket{1} & \bra{n}\hat{\Omega}\ket{2} & \cdots & \bra{n}\hat{\Omega}\ket{n}
    \end{bmatrix}
    \begin{bmatrix}
        v_1 \\
        v_2 \\
        \vdots \\
        v_n
    \end{bmatrix}
\end{equation*}

Multiplying two operators $\hat{\Omega}$ and $\hat{\Lambda}$ corresponds to multiplying their matrix representations, and can be interpreted as first applying $\hat{\Omega}$ to a vector and then applying $\hat{\Lambda}$ to the result. Thus, we have:
\begin{equation*}
    \hat{(\Lambda\Omega)} \ket{V} = \hat{\Lambda} (\hat{\Omega} \ket{V})
\end{equation*}

The order of operations is important. In general, $\hat{\Lambda}\hat{\Omega} \neq \hat{\Omega}\hat{\Lambda}$, and we define a new operator called the commutator of $\hat{\Omega}$ and $\hat{\Lambda}$ as:
\begin{equation*}
    [\hat{\Omega}, \hat{\Lambda}] = \hat{\Omega}\hat{\Lambda} - \hat{\Lambda}\hat{\Omega}
\end{equation*}

Some properties of the commutator are:
\begin{itemize}
    \item $[\hat{\Omega}, \hat{\Lambda}] = -[\hat{\Lambda}, \hat{\Omega}]$.
    \item $[\hat{\Omega}, \hat{\Lambda} + \hat{\Theta}] = [\hat{\Omega}, \hat{\Lambda}] + [\hat{\Omega}, \hat{\Theta}]$.
    \item $[\hat{\Omega} + \hat{\Theta}, \hat{\Lambda}] = [\hat{\Omega}, \hat{\Lambda}] + [\hat{\Theta}, \hat{\Lambda}]$.
    \item $[\hat{\Omega}, c\hat{\Lambda}] = c[\hat{\Omega}, \hat{\Lambda}]$, where $c$ is a scalar.
    \item $[c\hat{\Omega}, \hat{\Lambda}] = c[\hat{\Omega}, \hat{\Lambda}]$, where $c$ is a scalar.
    \item $[\hat{\Omega}, \hat{\Lambda}\hat{\Theta}] = [\hat{\Omega}, \hat{\Lambda}]\hat{\Theta} + \hat{\Lambda}[\hat{\Omega}, \hat{\Theta}]$.
    \item $[\hat{\Omega}\hat{\Lambda}, \hat{\Theta}] = \hat{\Omega}[\hat{\Lambda}, \hat{\Theta}] + [\hat{\Omega}, \hat{\Theta}]\hat{\Lambda}$.
\end{itemize}

The inverse of an operator $\hat{\Omega}$, if exists, is denoted as $\hat{\Omega}^{-1}$ and is the operator that satisfies the following conditions:
\begin{equation*}
    \hat{\Omega}^{-1} \hat{\Omega} = \hat{\Omega} \hat{\Omega}^{-1} = \hat{I}
\end{equation*}
\begin{equation*}
    \hat{(\Omega\Lambda)}^{-1} = \hat{\Lambda}^{-1} \hat{\Omega}^{-1}
\end{equation*}

\subsection{Adjoint of an Operator}

The adjoint of an operator $\hat{\Omega}$, denoted as $\hat{\Omega}^\dagger$, is the operator that satisfies the following condition for all vectors $\ket{V}$ and $\ket{W}$:

\begin{equation*}
    \bra{V}\hat{\Omega}\ket{W} = \bra{W}\hat{\Omega}^\dagger\ket{V}^*
\end{equation*}

Now, we have

\begin{equation*}
    \Omega^\dagger_{ij} = \bra{i}\hat{\Omega}^\dagger\ket{j} = \bra{j}\hat{\Omega}\ket{i}^* = \Omega_{ji}^*
\end{equation*}

Thus, the matrix representation of the adjoint operator $\hat{\Omega}^\dagger$ is the conjugate transpose of the matrix representation of $\hat{\Omega}$.


Some properties of the adjoint operation are:
\begin{itemize}
    \item $(\hat{\Omega} + \hat{\Lambda})^\dagger = \hat{\Omega}^\dagger + \hat{\Lambda}^\dagger$.
    \item $(c\hat{\Omega})^\dagger = \hat{\Omega}^\dagger c^*$, where $c$ is a scalar.
    \item $(\hat{\Omega}\hat{\Lambda})^\dagger = \hat{\Lambda}^\dagger \hat{\Omega}^\dagger$.
    \item $(\hat{\Omega}^\dagger)^\dagger = \hat{\Omega}$.
\end{itemize}

\subsection{Special Types of Operators}

\begin{itemize}

\item Hermitian and Anti-Hermitian Operators

An operator $\hat{\Omega}$ is called Hermitian if it is equal to its adjoint, i.e., $\hat{\Omega} = \hat{\Omega}^\dagger$.
An operator $\hat{\Omega}$ is called anti-Hermitian if it is equal to the negative of its adjoint, i.e., $\hat{\Omega} = -\hat{\Omega}^\dagger$.

Any operator $\hat{\Omega}$ can be expressed as a sum of a Hermitian operator and an anti-Hermitian operator:
\begin{equation*}
    \hat{\Omega} = \frac{1}{2}(\hat{\Omega} + \hat{\Omega}^\dagger) + \frac{1}{2}(\hat{\Omega} - \hat{\Omega}^\dagger)
\end{equation*}

\item Unitary Operators

An operator $\hat{U}$ is called unitary if it satisfies the following condition:
\begin{equation*}
    \hat{U}^\dagger \hat{U} = \hat{U} \hat{U}^\dagger = \hat{I}
\end{equation*}

Unitary operators represent transformations that preserve the inner product, and thus preserve the norm of vectors. They are often used to represent symmetries and time evolution in quantum mechanics.

To show that unitary operators preserve the inner product, we can calculate the inner product of two vectors $\ket{V}$ and $\ket{W}$ after applying a unitary operator $\hat{U}$:
\begin{equation*}
    \bra{\hat{U}V}\ket{\hat{U}W} = \bra{V}\hat{U}^\dagger \hat{U}\ket{W} = \bra{V}\hat{I}\ket{W} = \bra{V}\ket{W}
\end{equation*}

As a consequence, the individual columns of a unitary operator $\hat{U}$ form an orthonormal set of vectors. This is because the inner product of any two columns of $\hat{U}$ is given by:
\begin{equation*}
    \bra{\hat{U}i}\ket{\hat{U}j} = \bra{i}\hat{U}^\dagger \hat{U}\ket{j} = \bra{i}\hat{I}\ket{j} = \delta_{ij}
\end{equation*}

\end{itemize}

\subsection{Active and Passive Transformations}

Suppose all vectors in a vector space $\vecspace{V}$ is subjected to a unitary transformation $\hat{U}$.

\begin{equation*}
    \ket{V} \rightarrow \ket{V'} = \hat{U}\ket{V}
\end{equation*}

Under the transformation $\hat{U}$, the matrix representation of any operator $\hat{\Omega}$ is modified as:

\begin{equation*}
    \Omega'_{ij} = \bra{i'}\hat{\Omega}\ket{j'} = \bra{\hat{U}i}\hat{\Omega}\ket{\hat{U}j} = \bra{i}\hat{U}^\dagger\hat{\Omega}\hat{U}\ket{j} = (U^\dagger\Omega U)_{ij}
\end{equation*}

Thus, the change is similar to leaving the vectors alone and changing all operators to:

\begin{equation*}
    \hat{\Omega} \rightarrow \hat{U}^\dagger\hat{\Omega}\hat{U}
\end{equation*}

The first case is called an active transformation, while the second case is called a passive transformation.

\subsubsection{Trace of an Operator}

The trace of a matrix is defined as:

\begin{equation*}
    Tr(\Omega) = \sum_i \Omega_{ii}
\end{equation*}

Some properties of trace of an operator are:

\begin{itemize}
    \item $Tr(\Omega\Lambda) = Tr(\Lambda\Omega)$
    \item $Tr(\Omega\Lambda\Theta) = Tr(\Lambda\Theta\Omega) = Tr(\Theta\Omega\Lambda)$
    \item $Tr(\Omega) = Tr(U^\dagger\Omega U)$ where $U$ is a unitary matrix.
\end{itemize}


% -------------------------------------------------
% References
% -------------------------------------------------
\begin{thebibliography}{99}

\bibitem{griffiths}
David J. Griffiths,
\textit{Introduction to Quantum Mechanics},
3rd Edition.

\bibitem{sakurai}
J. J. Sakurai,
\textit{Modern Quantum Mechanics}.

\bibitem{shankar}
R. Shankar,
\textit{Principles of Quantum Mechanics}.

\end{thebibliography}

\end{document}